% ============================================================
%  bundle2.tex — Taxonomy, Literature Survey (5 domains),
%                Comparative Analysis
%  % ============================================================
%  bundle2.tex — Taxonomy, Literature Survey (5 domains),
%                Comparative Analysis
%  % ============================================================
%  bundle2.tex — Taxonomy, Literature Survey (5 domains),
%                Comparative Analysis
%  % ============================================================
%  bundle2.tex — Taxonomy, Literature Survey (5 domains),
%                Comparative Analysis
%  \input{sections/bundle2} from main.tex
% ============================================================

\section{Taxonomy of Related Work}
\label{sec:taxonomy}

The surveyed literature was organized into five functional domains that together cover the full processing pipeline of an integrated defence surveillance system: (A) multi-modal sensor fusion and data integration, (B) multi-agent AI for threat detection and classification, (C) false-alarm reduction and explainable decision support, (D) geospatial situational awareness and human-in-the-loop governance, and (E) real-time distributed architecture and adversarial robustness. To keep the scope focused and relevant, inclusion criteria limited the corpus to studies that satisfied at least two of three conditions: applicability to defence or security-critical environments, use of AI or machine learning methods, and involvement of multi-sensor or distributed system architectures. The sources drawn upon span peer-reviewed journal articles, IEEE and ACM conference proceedings, and policy-research technical reports~\cite{cset2024maven, cset2025ai}.

Table~\ref{tab:taxonomy} presents the resulting classification. For each domain, the table identifies representative published approaches, the principal sensor or data modalities examined, the key technique employed, and the primary limitation observed when that approach is evaluated against the demands of a unified multi-domain surveillance framework. It is worth noting that these five domains do not operate in isolation from one another in practice. The quality of sensor fusion at the data integration layer in Domain A directly shapes the classification confidence that downstream agents in Domain B can reliably achieve. Similarly, the time it takes to generate explanations in the decision-support layer of Domain C places real constraints on how quickly human commanders can complete their authorization loop in Domain D. These cross-domain dependencies are examined quantitatively in Section~V.



\section{Literature Survey}
\label{sec:survey}

\subsection{Multi-Modal Sensor Fusion and Data Integration}
\label{sec:fusion}

\subsubsection{Problem Definition}
One of the most fundamental challenges in defence surveillance is that the sensors deployed across a border do not speak the same language. Satellite imagery covers vast areas but only updates every few hours. UAV feeds capture narrow zones in high-frame-rate video. Ground-based sensors like seismic detectors, acoustic microphones, and infrared cameras produce continuous streams of numeric time-series data, each with its own noise characteristics and sampling rate~\cite{meng2023multisource}. Bringing all of these together into a single coherent, low-latency operational picture is the central technical challenge of sensor fusion for border security. India's diverse terrain makes this harder still, as the conditions across border sectors range from arid desert to dense riverine mangrove, and connectivity through radio and satellite links is frequently intermittent, adding further complexity to any integration effort.

\subsubsection{Existing Approaches}
The most widely referenced framework for structuring this problem is the Joint Directors of Laboratories (JDL) data fusion model, which defines five processing levels ranging from raw signal refinement at Level~0 through situation assessment at Level~2 and up to threat refinement at Level~3. It remains the dominant architectural reference for military fusion systems worldwide~\cite{meng2023multisource}. A comprehensive survey of multi-source information fusion for command, control, communications, and intelligence systems by Meng et al.~\cite{meng2023multisource} draws a useful distinction between data-level, feature-level, and decision-level fusion, mapping out the computational tradeoffs involved at each stage. In counter-UAV applications where individual sensor reliability varies across optical, thermal, radio-frequency, and radar inputs, attention-based cross-modal feature alignment has consistently produced the strongest classification accuracy. At the industry level, defence contractors have pushed this further by integrating hyperspectral imaging with AI-enabled radar and LiDAR processing for military situational awareness, and passive infrared sensor networks fused with multi-modal inputs have already been deployed for autonomous detect-classify-track operations at military installations. Ali et al.~\cite{threatdetect2024rt} survey transformer-based models and federated learning pipelines for real-time threat detection, identifying attention mechanisms as particularly effective for aligning features across sensor modalities.

\subsubsection{Key Insight and Research Gap}
Despite this progress, a consistent limitation runs through the reviewed literature. Nearly all fusion architectures have been evaluated on homogeneous or pairwise sensor combinations and have not been tested against the five-or-more modality configurations that are actually present in deployed systems like India's Comprehensive Integrated Border Management System (CIBMS). The temporal misalignment between high-frequency ground sensors reporting at 100~Hz or above and low-frequency satellite revisit passes separated by hours remains a genuinely unsolved synchronisation problem~\cite{meng2023multisource}. Perhaps most critically, not a single study in the reviewed corpus evaluates fusion performance under adversarial sensor degradation conditions. Given the documented electromagnetic interference threats present in contested border environments~\cite{milai2025risk}, this is not a minor academic gap but a direct operational vulnerability that future work urgently needs to address.


\subsection{Multi-Agent AI for Threat Detection and Classification}
\label{sec:multiagent}

\subsubsection{Problem Definition}
A single AI model simply cannot do everything at once. Asking one model to simultaneously process live video from dozens of cameras, classify radar anomalies, analyse seismic and acoustic time-series data, and correlate all of it into a unified threat picture within the sub-second response windows that operational protocols demand is not realistic~\cite{zhao2022deepai}. The natural solution is to distribute these responsibilities across multiple specialised agents, each focused on what it does best. But this introduces its own set of challenges. These agents need to share findings with each other, reconcile situations where two agents reach contradictory conclusions about the same event, and ultimately agree on a single threat classification without any one agent becoming a bottleneck that slows the entire system down~\cite{hassan2025multiagent}.

\subsubsection{Existing Approaches}
Several recent systems have explored how to structure this kind of multi-agent design effectively. Hassan et al.~\cite{hassan2025multiagent} propose a modular multi-agent system for cybersecurity threat detection where separate agents handle log analysis, IP analysis, and threat correlation independently before feeding into a central correlation engine. Their results show that this role-based decomposition improves both detection accuracy and the ability to trace why a particular classification was made, compared to a single monolithic model. Zhao et al.~\cite{zhao2022deepai} describe a cooperative AI staff system managing over 500{,}000 simulated battlefield agents, delivering command-level situational awareness through automated spatial and temporal pattern correlation across a large operational area. The AgenticCyber framework~\cite{agenticcyber2025} takes multi-agent orchestration further by fusing cloud logs, surveillance video, and audio streams across modalities, reporting a 96.2\% F1 score at 420~milliseconds end-to-end latency. On the visual detection side, YOLOv8-family architectures have been successfully adapted for military targets in satellite imagery~\cite{satellite2024yolov8} and weapon identification in live surveillance feeds~\cite{weapon2023yolov8}. Project Maven's real-world deployment offered perhaps the most striking validation of the approach, with a 20-person AI-assisted cell matching the video analysis throughput of a 2{,}000-person manual team~\cite{cset2024maven}, though independent evaluation found tank identification accuracy was only around 60\% even under favourable conditions~\cite{rand2023maven}. Frameworks for formally dividing decision authority between AI agents and human commanders in tactical settings have also been developed~\cite{autonomyteaming2024}.

\subsubsection{Key Insight and Research Gap}
Looking across these systems, a clear and important gap emerges. Existing multi-agent surveillance architectures have largely been designed around uniform task distributions, such as passing a tracked target between cameras in a consistent video network. None of them seriously tackle the problem of orchestrating agents that are reasoning over fundamentally different types of data at the same time, where one agent is working with pixel arrays, another with audio spectrograms, and another with radar range-Doppler matrices~\cite{hassan2025multiagent, agenticcyber2025}. This matters enormously for a system like VIGILANCE. A graduated five-level threat classification scheme requires every individual agent to produce well-calibrated confidence scores, and the correlation agent must then perform principled uncertainty quantification when combining outputs from agents that each have their own distinct error profiles and failure modes. In current multi-agent designs, confidence calibration is treated as a secondary concern rather than something built into the architecture from the ground up~\cite{cset2025ai}. That is a gap this work directly sets out to address.


\subsection{False-Alarm Reduction and Explainable Decision Support}
\label{sec:falsealarm}

\subsubsection{Problem Definition}
Persistently high false-alarm rates constitute a documented and operationally consequential failure mode in fielded surveillance systems~\cite{deepsurvey2023anomaly}. Sustained spurious alerts induce operator alert fatigue, misallocate rapid-response resources, and progressively degrade trust in automated detection pipelines. Root causes include environmental noise from wind-induced fence vibration, animal crossings, and weather-driven radar clutter; absence of temporal consistency verification in threshold-based alarm logic; and lack of contextual priors that distinguish known civilian activity zones from restricted perimeters~\cite{cctv2024anomaly}.

\subsubsection{Existing Approaches}
Weakly-supervised multiple-instance learning (MIL) frameworks for anomaly detection in real-world surveillance video have established training paradigms that substantially reduce annotation requirements for large-scale deployment~\cite{deepsurvey2023anomaly}. Deep-learning surveys catalogue CNN, LSTM, and transformer-based temporal models that improve anomaly discrimination by enforcing signature persistence across consecutive frames before triggering an alert. Hybrid architectures combining spatial object detection with temporal verification have been shown to reduce false-positive rates relative to single-frame classifiers in CCTV monitoring~\cite{cctv2024anomaly}. On the explainability front, Zhang et al.~\cite{zhang2022xai} survey SHAP, LIME, and attention-based explanation mechanisms applicable to black-box classifiers, applying explainable AI (XAI) to intrusion detection and malware classification, tasks that are structurally analogous to the threat-triage function required in military surveillance. Roff and Danks~\cite{military2024xai} address military-specific accountability requirements, stipulating that commanders must be able to audit the rationale behind targeting and escalatory recommendations under both domestic and international legal obligations.

\subsubsection{Key Insight and Research Gap}
XAI methods in surveillance have been evaluated primarily for post-hoc interpretability on static or pre-recorded datasets~\cite{zhang2022xai} and have not been integrated into real-time decision pipelines where explanation generation latency itself constrains operator response time. Contextual filtering approaches that suppress alarms in geofenced civilian zones depend on accurate, current civilian activity maps that are unavailable or classified in active border areas~\cite{cctv2024anomaly}, creating an unresolved data-dependency bottleneck that no reviewed study has addressed.


\subsection{Geospatial Situational Awareness and Human-in-the-Loop Systems}
\label{sec:geospatial}

\subsubsection{Problem Definition}
Actionable command decisions require threat intelligence to be spatially anchored: operators must perceive not only the nature and severity of a detected anomaly but also its precise location relative to terrain features, installation boundaries, sensor coverage envelopes, and proximity to civilian population centres~\cite{karimipour2021iot}. Graduated threat classification further demands that alerts assessed at the highest severity levels do not trigger autonomous system action; human authorisation must remain mandatory at those tiers~\cite{meaningful2024hitl}. The interface design challenge is twofold: preventing operators from reflexively endorsing AI recommendations while preserving response latencies compatible with tactical decision cycles~\cite{hitl2023trust}.

\subsubsection{Existing Approaches}
Karimipour and Derakhshan~\cite{karimipour2021iot} describe AI-enabled situational awareness dashboards for distributed sensor networks in industrial Internet of Things (IoT) environments, incorporating geographic information system (GIS) threat mapping with sensor registries that visualise coverage zones, detection radii, and active alarm sources on unified terrain layers. Zhao et al.~\cite{zhao2022deepai} extend this concept to military command by constructing multi-layer spatio-temporal battlefield visualisations that overlay force positions, predicted adversary trajectories, and terrain accessibility indices. Geospatial UAV detection with movement-vector plotting on georeferenced maps has been demonstrated for counter-drone applications. Maas~\cite{hitl2023trust} examines trust calibration and automation bias in military human-autonomy teaming, reporting that operator trust is sensitive to the transparency and perceived reliability of AI outputs. Svenmarck and Luotsinen~\cite{meaningful2024hitl} formalise the legal and technical requirements for meaningful human control, proposing design constraints that ensure commander authority over engagement decisions is not diluted by automated recommendation systems.

\subsubsection{Key Insight and Research Gap}
Human-in-the-loop (HITL) frameworks in the surveyed literature focus on single-operator, single-threat interaction models~\cite{hitl2023trust, meaningful2024hitl}. Concurrent multi-threat environments, in which an operator must triage AI-generated alerts across geographically distributed incidents simultaneously, remain under-examined despite constituting the operational norm in border-surveillance command centres. Interface designs that communicate prediction uncertainty quantitatively to operators with varying levels of technical expertise present an open challenge; existing prototypes either display raw probability values that non-technical operators misinterpret or suppress uncertainty information entirely, promoting uncalibrated trust~\cite{hitl2023trust}.


\subsection{Real-Time Distributed Architecture and Adversarial Robustness}
\label{sec:distributed}

\subsubsection{Problem Definition}
Sensor data generated at a remote border outpost must traverse constrained communication links, undergo inference at sub-second latency, and arrive at a central command dashboard with actionable fidelity~\cite{flink2024kafka}. Meeting this requirement demands edge preprocessing to reduce transmission volume, stream-processing middleware for sequential low-latency inference, and an architecture tolerant of node failures and intermittent connectivity. Adversarial robustness introduces a distinct threat surface: a high-value surveillance system constitutes an explicit target for adversarial patches designed to defeat object detectors, GPS spoofing aimed at unmanned platform navigation, and signal jamming directed at ground-sensor communication uplinks~\cite{advyolo2022, milai2025risk}.

\subsubsection{Existing Approaches}
Patel et al.~\cite{flink2024kafka} present a real-time IoT sensor streaming architecture combining Apache Kafka for telemetry ingestion with Apache Flink for stream processing, achieving low-latency anomaly detection on continuous sensor feeds with fault-tolerant horizontal scaling. Edge-computing extensions of this pattern perform preliminary inference at the sensor node before transmitting compressed results, reducing bandwidth requirements for drone and CCTV data in remote deployments. Ali et al.~\cite{threatdetect2024rt} survey transformer-based models and federated learning pipelines for real-time threat detection across IoT and 5G networks, identifying federated architectures as particularly suitable for privacy-preserving distributed inference in military contexts. On the adversarial front, Liang et al.~\cite{advyolo2022} catalogue physical-world attacks on YOLO-family detectors, including adversarial patches and printed perturbations that cause vehicle misclassification in outdoor conditions. Wang et al.~\cite{milai2025risk} address military-specific adversarial risks such as electromagnetic interference and GPS spoofing, and propose ethical governance guardrails for deployed military AI. Scientometric analysis of approximately 3{,}500 Scopus records confirms a fourfold growth in AI-related defence publications between 2015 and 2025~\cite{emerging2025scientometric}, yet distributed surveillance architectures remain under-represented relative to single-algorithm detection studies. Independent PRISMA reviews of UAV cybersecurity~\cite{uavsec2025iet} and UAV intrusion detection systems~\cite{uavids2025drones} converge on a shared conclusion: no unified cross-domain framework currently addresses the full defence surveillance pipeline from edge to command centre.

\subsubsection{Key Insight and Research Gap}
Adversarial robustness research in computer vision predominantly addresses digital perturbations on laboratory-curated datasets; physical-world attacks under outdoor surveillance conditions involving variable lighting, partial occlusion, and ambient sensor noise are substantially less investigated~\cite{advyolo2022}. No reviewed distributed-architecture study evaluates end-to-end pipeline latency under simultaneous edge-node failure and adversarial sensor injection, the compound threat model most representative of a contested border environment. The gap between adversarial machine-learning research and operational distributed systems engineering constitutes one of the least-addressed intersections in the defence surveillance literature~\cite{uavsec2025iet, milai2025risk}.


\section{Comparative Analysis}
\label{sec:comparative}

Cross-domain evaluation reveals an uneven maturity gradient across the five surveyed areas. Sensor fusion (Domain~A) and visual object detection (a component of Domain~B) represent the most technically mature pillars, and YOLOv8-family detectors~\cite{satellite2024yolov8, weapon2023yolov8} approach deployment readiness for constrained single-modality scenarios. False-alarm reduction techniques (Domain~C) yield measurable improvements on benchmark datasets~\cite{deepsurvey2023anomaly} but lack validated integration with real-time explainability pipelines and with the accountability demands specific to military decision chains~\cite{military2024xai}. Geospatial HITL systems (Domain~D) and adversarially hardened distributed architectures (Domain~E) remain at the proof-of-concept stage, with empirical validation limited to single-site deployments or simulated environments~\cite{karimipour2021iot, flink2024kafka}.

The most consequential pattern emerging from this review is the near-total absence of research addressing interaction effects across domain boundaries. Fusion fidelity at the sensor-integration layer directly determines the confidence bounds available to downstream classification agents~\cite{meng2023multisource, hassan2025multiagent}. Explanation-generation overhead in the XAI layer imposes an additive latency penalty on the HITL authorisation cycle~\cite{zhang2022xai, hitl2023trust}. Adversarial corruption at edge sensors propagates cascading failures through fusion, classification, and alerting stages that no surveyed system models as an integrated failure mode~\cite{advyolo2022, milai2025risk}. Table~\ref{tab:comparative} summarises the best-performing technique, reported benchmark, metric, and principal limitation for each domain when assessed against the requirements of a unified multi-domain surveillance framework.

\begin{table*}[!t]
\centering
\caption{Comparative Analysis of Surveyed Approaches}
\label{tab:comparative}
\renewcommand{\arraystretch}{1.3}
\footnotesize
\begin{tabular}{>{\raggedright\arraybackslash}p{1.5cm}>{\raggedright\arraybackslash}p{3.2cm}>{\raggedright\arraybackslash}p{3.0cm}>{\raggedright\arraybackslash}p{2.8cm}>{\raggedright\arraybackslash}p{3.5cm}}
\hline
\textbf{Domain} & \textbf{Best Performing Technique} & \textbf{Dataset / Benchmark} & \textbf{Metric Reported} & \textbf{Limitation in Unified Context} \\
\hline
A: Sensor Fusion
  & Attention-based cross-modal DL fusion
  & Counter-UAV optical, thermal, RF, radar corpus
  & Highest classification accuracy vs.\ single-modal baselines
  & No temporal synchronisation for $>$2 modalities at disparate sampling rates \\
\hline
B: Multi-Agent AI
  & Multi-modal agentic orchestration
  & Cloud logs + surveillance video + audio
  & 96.2\% F1; 420\,ms end-to-end latency
  & Agent confidence calibration treated as post-hoc addition \\
\hline
C: False Alarm / XAI
  & Deep-learning anomaly detection with temporal verification
  & Real-world surveillance video benchmarks
  & Benchmark AUC across anomaly classes
  & Explanation latency not measured for real-time deployment \\
\hline
D: Geospatial / HITL
  & Multi-layer battlefield visualisation
  & 500K+ simulated battlefield agents
  & Command-level situational awareness coverage
  & Multi-threat concurrent operator triage not evaluated \\
\hline
E: Distributed / Adversarial
  & Kafka + Flink real-time streaming
  & Continuous IoT sensor streams
  & Low-latency, fault-tolerant stream processing
  & Not tested under adversarial sensor injection or node compromise \\
\hline
\end{tabular}
\end{table*}
 from main.tex
% ============================================================

\section{Taxonomy of Related Work}
\label{sec:taxonomy}

The surveyed literature was organized into five functional domains that together cover the full processing pipeline of an integrated defence surveillance system: (A) multi-modal sensor fusion and data integration, (B) multi-agent AI for threat detection and classification, (C) false-alarm reduction and explainable decision support, (D) geospatial situational awareness and human-in-the-loop governance, and (E) real-time distributed architecture and adversarial robustness. To keep the scope focused and relevant, inclusion criteria limited the corpus to studies that satisfied at least two of three conditions: applicability to defence or security-critical environments, use of AI or machine learning methods, and involvement of multi-sensor or distributed system architectures. The sources drawn upon span peer-reviewed journal articles, IEEE and ACM conference proceedings, and policy-research technical reports~\cite{cset2024maven, cset2025ai}.

Table~\ref{tab:taxonomy} presents the resulting classification. For each domain, the table identifies representative published approaches, the principal sensor or data modalities examined, the key technique employed, and the primary limitation observed when that approach is evaluated against the demands of a unified multi-domain surveillance framework. It is worth noting that these five domains do not operate in isolation from one another in practice. The quality of sensor fusion at the data integration layer in Domain A directly shapes the classification confidence that downstream agents in Domain B can reliably achieve. Similarly, the time it takes to generate explanations in the decision-support layer of Domain C places real constraints on how quickly human commanders can complete their authorization loop in Domain D. These cross-domain dependencies are examined quantitatively in Section~V.



\section{Literature Survey}
\label{sec:survey}

\subsection{Multi-Modal Sensor Fusion and Data Integration}
\label{sec:fusion}

\subsubsection{Problem Definition}
One of the most fundamental challenges in defence surveillance is that the sensors deployed across a border do not speak the same language. Satellite imagery covers vast areas but only updates every few hours. UAV feeds capture narrow zones in high-frame-rate video. Ground-based sensors like seismic detectors, acoustic microphones, and infrared cameras produce continuous streams of numeric time-series data, each with its own noise characteristics and sampling rate~\cite{meng2023multisource}. Bringing all of these together into a single coherent, low-latency operational picture is the central technical challenge of sensor fusion for border security. India's diverse terrain makes this harder still, as the conditions across border sectors range from arid desert to dense riverine mangrove, and connectivity through radio and satellite links is frequently intermittent, adding further complexity to any integration effort.

\subsubsection{Existing Approaches}
The most widely referenced framework for structuring this problem is the Joint Directors of Laboratories (JDL) data fusion model, which defines five processing levels ranging from raw signal refinement at Level~0 through situation assessment at Level~2 and up to threat refinement at Level~3. It remains the dominant architectural reference for military fusion systems worldwide~\cite{meng2023multisource}. A comprehensive survey of multi-source information fusion for command, control, communications, and intelligence systems by Meng et al.~\cite{meng2023multisource} draws a useful distinction between data-level, feature-level, and decision-level fusion, mapping out the computational tradeoffs involved at each stage. In counter-UAV applications where individual sensor reliability varies across optical, thermal, radio-frequency, and radar inputs, attention-based cross-modal feature alignment has consistently produced the strongest classification accuracy. At the industry level, defence contractors have pushed this further by integrating hyperspectral imaging with AI-enabled radar and LiDAR processing for military situational awareness, and passive infrared sensor networks fused with multi-modal inputs have already been deployed for autonomous detect-classify-track operations at military installations. Ali et al.~\cite{threatdetect2024rt} survey transformer-based models and federated learning pipelines for real-time threat detection, identifying attention mechanisms as particularly effective for aligning features across sensor modalities.

\subsubsection{Key Insight and Research Gap}
Despite this progress, a consistent limitation runs through the reviewed literature. Nearly all fusion architectures have been evaluated on homogeneous or pairwise sensor combinations and have not been tested against the five-or-more modality configurations that are actually present in deployed systems like India's Comprehensive Integrated Border Management System (CIBMS). The temporal misalignment between high-frequency ground sensors reporting at 100~Hz or above and low-frequency satellite revisit passes separated by hours remains a genuinely unsolved synchronisation problem~\cite{meng2023multisource}. Perhaps most critically, not a single study in the reviewed corpus evaluates fusion performance under adversarial sensor degradation conditions. Given the documented electromagnetic interference threats present in contested border environments~\cite{milai2025risk}, this is not a minor academic gap but a direct operational vulnerability that future work urgently needs to address.


\subsection{Multi-Agent AI for Threat Detection and Classification}
\label{sec:multiagent}

\subsubsection{Problem Definition}
A single AI model simply cannot do everything at once. Asking one model to simultaneously process live video from dozens of cameras, classify radar anomalies, analyse seismic and acoustic time-series data, and correlate all of it into a unified threat picture within the sub-second response windows that operational protocols demand is not realistic~\cite{zhao2022deepai}. The natural solution is to distribute these responsibilities across multiple specialised agents, each focused on what it does best. But this introduces its own set of challenges. These agents need to share findings with each other, reconcile situations where two agents reach contradictory conclusions about the same event, and ultimately agree on a single threat classification without any one agent becoming a bottleneck that slows the entire system down~\cite{hassan2025multiagent}.

\subsubsection{Existing Approaches}
Several recent systems have explored how to structure this kind of multi-agent design effectively. Hassan et al.~\cite{hassan2025multiagent} propose a modular multi-agent system for cybersecurity threat detection where separate agents handle log analysis, IP analysis, and threat correlation independently before feeding into a central correlation engine. Their results show that this role-based decomposition improves both detection accuracy and the ability to trace why a particular classification was made, compared to a single monolithic model. Zhao et al.~\cite{zhao2022deepai} describe a cooperative AI staff system managing over 500{,}000 simulated battlefield agents, delivering command-level situational awareness through automated spatial and temporal pattern correlation across a large operational area. The AgenticCyber framework~\cite{agenticcyber2025} takes multi-agent orchestration further by fusing cloud logs, surveillance video, and audio streams across modalities, reporting a 96.2\% F1 score at 420~milliseconds end-to-end latency. On the visual detection side, YOLOv8-family architectures have been successfully adapted for military targets in satellite imagery~\cite{satellite2024yolov8} and weapon identification in live surveillance feeds~\cite{weapon2023yolov8}. Project Maven's real-world deployment offered perhaps the most striking validation of the approach, with a 20-person AI-assisted cell matching the video analysis throughput of a 2{,}000-person manual team~\cite{cset2024maven}, though independent evaluation found tank identification accuracy was only around 60\% even under favourable conditions~\cite{rand2023maven}. Frameworks for formally dividing decision authority between AI agents and human commanders in tactical settings have also been developed~\cite{autonomyteaming2024}.

\subsubsection{Key Insight and Research Gap}
Looking across these systems, a clear and important gap emerges. Existing multi-agent surveillance architectures have largely been designed around uniform task distributions, such as passing a tracked target between cameras in a consistent video network. None of them seriously tackle the problem of orchestrating agents that are reasoning over fundamentally different types of data at the same time, where one agent is working with pixel arrays, another with audio spectrograms, and another with radar range-Doppler matrices~\cite{hassan2025multiagent, agenticcyber2025}. This matters enormously for a system like VIGILANCE. A graduated five-level threat classification scheme requires every individual agent to produce well-calibrated confidence scores, and the correlation agent must then perform principled uncertainty quantification when combining outputs from agents that each have their own distinct error profiles and failure modes. In current multi-agent designs, confidence calibration is treated as a secondary concern rather than something built into the architecture from the ground up~\cite{cset2025ai}. That is a gap this work directly sets out to address.


\subsection{False-Alarm Reduction and Explainable Decision Support}
\label{sec:falsealarm}

\subsubsection{Problem Definition}
Persistently high false-alarm rates constitute a documented and operationally consequential failure mode in fielded surveillance systems~\cite{deepsurvey2023anomaly}. Sustained spurious alerts induce operator alert fatigue, misallocate rapid-response resources, and progressively degrade trust in automated detection pipelines. Root causes include environmental noise from wind-induced fence vibration, animal crossings, and weather-driven radar clutter; absence of temporal consistency verification in threshold-based alarm logic; and lack of contextual priors that distinguish known civilian activity zones from restricted perimeters~\cite{cctv2024anomaly}.

\subsubsection{Existing Approaches}
Weakly-supervised multiple-instance learning (MIL) frameworks for anomaly detection in real-world surveillance video have established training paradigms that substantially reduce annotation requirements for large-scale deployment~\cite{deepsurvey2023anomaly}. Deep-learning surveys catalogue CNN, LSTM, and transformer-based temporal models that improve anomaly discrimination by enforcing signature persistence across consecutive frames before triggering an alert. Hybrid architectures combining spatial object detection with temporal verification have been shown to reduce false-positive rates relative to single-frame classifiers in CCTV monitoring~\cite{cctv2024anomaly}. On the explainability front, Zhang et al.~\cite{zhang2022xai} survey SHAP, LIME, and attention-based explanation mechanisms applicable to black-box classifiers, applying explainable AI (XAI) to intrusion detection and malware classification, tasks that are structurally analogous to the threat-triage function required in military surveillance. Roff and Danks~\cite{military2024xai} address military-specific accountability requirements, stipulating that commanders must be able to audit the rationale behind targeting and escalatory recommendations under both domestic and international legal obligations.

\subsubsection{Key Insight and Research Gap}
XAI methods in surveillance have been evaluated primarily for post-hoc interpretability on static or pre-recorded datasets~\cite{zhang2022xai} and have not been integrated into real-time decision pipelines where explanation generation latency itself constrains operator response time. Contextual filtering approaches that suppress alarms in geofenced civilian zones depend on accurate, current civilian activity maps that are unavailable or classified in active border areas~\cite{cctv2024anomaly}, creating an unresolved data-dependency bottleneck that no reviewed study has addressed.


\subsection{Geospatial Situational Awareness and Human-in-the-Loop Systems}
\label{sec:geospatial}

\subsubsection{Problem Definition}
Actionable command decisions require threat intelligence to be spatially anchored: operators must perceive not only the nature and severity of a detected anomaly but also its precise location relative to terrain features, installation boundaries, sensor coverage envelopes, and proximity to civilian population centres~\cite{karimipour2021iot}. Graduated threat classification further demands that alerts assessed at the highest severity levels do not trigger autonomous system action; human authorisation must remain mandatory at those tiers~\cite{meaningful2024hitl}. The interface design challenge is twofold: preventing operators from reflexively endorsing AI recommendations while preserving response latencies compatible with tactical decision cycles~\cite{hitl2023trust}.

\subsubsection{Existing Approaches}
Karimipour and Derakhshan~\cite{karimipour2021iot} describe AI-enabled situational awareness dashboards for distributed sensor networks in industrial Internet of Things (IoT) environments, incorporating geographic information system (GIS) threat mapping with sensor registries that visualise coverage zones, detection radii, and active alarm sources on unified terrain layers. Zhao et al.~\cite{zhao2022deepai} extend this concept to military command by constructing multi-layer spatio-temporal battlefield visualisations that overlay force positions, predicted adversary trajectories, and terrain accessibility indices. Geospatial UAV detection with movement-vector plotting on georeferenced maps has been demonstrated for counter-drone applications. Maas~\cite{hitl2023trust} examines trust calibration and automation bias in military human-autonomy teaming, reporting that operator trust is sensitive to the transparency and perceived reliability of AI outputs. Svenmarck and Luotsinen~\cite{meaningful2024hitl} formalise the legal and technical requirements for meaningful human control, proposing design constraints that ensure commander authority over engagement decisions is not diluted by automated recommendation systems.

\subsubsection{Key Insight and Research Gap}
Human-in-the-loop (HITL) frameworks in the surveyed literature focus on single-operator, single-threat interaction models~\cite{hitl2023trust, meaningful2024hitl}. Concurrent multi-threat environments, in which an operator must triage AI-generated alerts across geographically distributed incidents simultaneously, remain under-examined despite constituting the operational norm in border-surveillance command centres. Interface designs that communicate prediction uncertainty quantitatively to operators with varying levels of technical expertise present an open challenge; existing prototypes either display raw probability values that non-technical operators misinterpret or suppress uncertainty information entirely, promoting uncalibrated trust~\cite{hitl2023trust}.


\subsection{Real-Time Distributed Architecture and Adversarial Robustness}
\label{sec:distributed}

\subsubsection{Problem Definition}
Sensor data generated at a remote border outpost must traverse constrained communication links, undergo inference at sub-second latency, and arrive at a central command dashboard with actionable fidelity~\cite{flink2024kafka}. Meeting this requirement demands edge preprocessing to reduce transmission volume, stream-processing middleware for sequential low-latency inference, and an architecture tolerant of node failures and intermittent connectivity. Adversarial robustness introduces a distinct threat surface: a high-value surveillance system constitutes an explicit target for adversarial patches designed to defeat object detectors, GPS spoofing aimed at unmanned platform navigation, and signal jamming directed at ground-sensor communication uplinks~\cite{advyolo2022, milai2025risk}.

\subsubsection{Existing Approaches}
Patel et al.~\cite{flink2024kafka} present a real-time IoT sensor streaming architecture combining Apache Kafka for telemetry ingestion with Apache Flink for stream processing, achieving low-latency anomaly detection on continuous sensor feeds with fault-tolerant horizontal scaling. Edge-computing extensions of this pattern perform preliminary inference at the sensor node before transmitting compressed results, reducing bandwidth requirements for drone and CCTV data in remote deployments. Ali et al.~\cite{threatdetect2024rt} survey transformer-based models and federated learning pipelines for real-time threat detection across IoT and 5G networks, identifying federated architectures as particularly suitable for privacy-preserving distributed inference in military contexts. On the adversarial front, Liang et al.~\cite{advyolo2022} catalogue physical-world attacks on YOLO-family detectors, including adversarial patches and printed perturbations that cause vehicle misclassification in outdoor conditions. Wang et al.~\cite{milai2025risk} address military-specific adversarial risks such as electromagnetic interference and GPS spoofing, and propose ethical governance guardrails for deployed military AI. Scientometric analysis of approximately 3{,}500 Scopus records confirms a fourfold growth in AI-related defence publications between 2015 and 2025~\cite{emerging2025scientometric}, yet distributed surveillance architectures remain under-represented relative to single-algorithm detection studies. Independent PRISMA reviews of UAV cybersecurity~\cite{uavsec2025iet} and UAV intrusion detection systems~\cite{uavids2025drones} converge on a shared conclusion: no unified cross-domain framework currently addresses the full defence surveillance pipeline from edge to command centre.

\subsubsection{Key Insight and Research Gap}
Adversarial robustness research in computer vision predominantly addresses digital perturbations on laboratory-curated datasets; physical-world attacks under outdoor surveillance conditions involving variable lighting, partial occlusion, and ambient sensor noise are substantially less investigated~\cite{advyolo2022}. No reviewed distributed-architecture study evaluates end-to-end pipeline latency under simultaneous edge-node failure and adversarial sensor injection, the compound threat model most representative of a contested border environment. The gap between adversarial machine-learning research and operational distributed systems engineering constitutes one of the least-addressed intersections in the defence surveillance literature~\cite{uavsec2025iet, milai2025risk}.


\section{Comparative Analysis}
\label{sec:comparative}

Cross-domain evaluation reveals an uneven maturity gradient across the five surveyed areas. Sensor fusion (Domain~A) and visual object detection (a component of Domain~B) represent the most technically mature pillars, and YOLOv8-family detectors~\cite{satellite2024yolov8, weapon2023yolov8} approach deployment readiness for constrained single-modality scenarios. False-alarm reduction techniques (Domain~C) yield measurable improvements on benchmark datasets~\cite{deepsurvey2023anomaly} but lack validated integration with real-time explainability pipelines and with the accountability demands specific to military decision chains~\cite{military2024xai}. Geospatial HITL systems (Domain~D) and adversarially hardened distributed architectures (Domain~E) remain at the proof-of-concept stage, with empirical validation limited to single-site deployments or simulated environments~\cite{karimipour2021iot, flink2024kafka}.

The most consequential pattern emerging from this review is the near-total absence of research addressing interaction effects across domain boundaries. Fusion fidelity at the sensor-integration layer directly determines the confidence bounds available to downstream classification agents~\cite{meng2023multisource, hassan2025multiagent}. Explanation-generation overhead in the XAI layer imposes an additive latency penalty on the HITL authorisation cycle~\cite{zhang2022xai, hitl2023trust}. Adversarial corruption at edge sensors propagates cascading failures through fusion, classification, and alerting stages that no surveyed system models as an integrated failure mode~\cite{advyolo2022, milai2025risk}. Table~\ref{tab:comparative} summarises the best-performing technique, reported benchmark, metric, and principal limitation for each domain when assessed against the requirements of a unified multi-domain surveillance framework.

\begin{table*}[!t]
\centering
\caption{Comparative Analysis of Surveyed Approaches}
\label{tab:comparative}
\renewcommand{\arraystretch}{1.3}
\footnotesize
\begin{tabular}{>{\raggedright\arraybackslash}p{1.5cm}>{\raggedright\arraybackslash}p{3.2cm}>{\raggedright\arraybackslash}p{3.0cm}>{\raggedright\arraybackslash}p{2.8cm}>{\raggedright\arraybackslash}p{3.5cm}}
\hline
\textbf{Domain} & \textbf{Best Performing Technique} & \textbf{Dataset / Benchmark} & \textbf{Metric Reported} & \textbf{Limitation in Unified Context} \\
\hline
A: Sensor Fusion
  & Attention-based cross-modal DL fusion
  & Counter-UAV optical, thermal, RF, radar corpus
  & Highest classification accuracy vs.\ single-modal baselines
  & No temporal synchronisation for $>$2 modalities at disparate sampling rates \\
\hline
B: Multi-Agent AI
  & Multi-modal agentic orchestration
  & Cloud logs + surveillance video + audio
  & 96.2\% F1; 420\,ms end-to-end latency
  & Agent confidence calibration treated as post-hoc addition \\
\hline
C: False Alarm / XAI
  & Deep-learning anomaly detection with temporal verification
  & Real-world surveillance video benchmarks
  & Benchmark AUC across anomaly classes
  & Explanation latency not measured for real-time deployment \\
\hline
D: Geospatial / HITL
  & Multi-layer battlefield visualisation
  & 500K+ simulated battlefield agents
  & Command-level situational awareness coverage
  & Multi-threat concurrent operator triage not evaluated \\
\hline
E: Distributed / Adversarial
  & Kafka + Flink real-time streaming
  & Continuous IoT sensor streams
  & Low-latency, fault-tolerant stream processing
  & Not tested under adversarial sensor injection or node compromise \\
\hline
\end{tabular}
\end{table*}
 from main.tex
% ============================================================

\section{Taxonomy of Related Work}
\label{sec:taxonomy}

The surveyed literature was organized into five functional domains that together cover the full processing pipeline of an integrated defence surveillance system: (A) multi-modal sensor fusion and data integration, (B) multi-agent AI for threat detection and classification, (C) false-alarm reduction and explainable decision support, (D) geospatial situational awareness and human-in-the-loop governance, and (E) real-time distributed architecture and adversarial robustness. To keep the scope focused and relevant, inclusion criteria limited the corpus to studies that satisfied at least two of three conditions: applicability to defence or security-critical environments, use of AI or machine learning methods, and involvement of multi-sensor or distributed system architectures. The sources drawn upon span peer-reviewed journal articles, IEEE and ACM conference proceedings, and policy-research technical reports~\cite{cset2024maven, cset2025ai}.

Table~\ref{tab:taxonomy} presents the resulting classification. For each domain, the table identifies representative published approaches, the principal sensor or data modalities examined, the key technique employed, and the primary limitation observed when that approach is evaluated against the demands of a unified multi-domain surveillance framework. It is worth noting that these five domains do not operate in isolation from one another in practice. The quality of sensor fusion at the data integration layer in Domain A directly shapes the classification confidence that downstream agents in Domain B can reliably achieve. Similarly, the time it takes to generate explanations in the decision-support layer of Domain C places real constraints on how quickly human commanders can complete their authorization loop in Domain D. These cross-domain dependencies are examined quantitatively in Section~V.



\section{Literature Survey}
\label{sec:survey}

\subsection{Multi-Modal Sensor Fusion and Data Integration}
\label{sec:fusion}

\subsubsection{Problem Definition}
One of the most fundamental challenges in defence surveillance is that the sensors deployed across a border do not speak the same language. Satellite imagery covers vast areas but only updates every few hours. UAV feeds capture narrow zones in high-frame-rate video. Ground-based sensors like seismic detectors, acoustic microphones, and infrared cameras produce continuous streams of numeric time-series data, each with its own noise characteristics and sampling rate~\cite{meng2023multisource}. Bringing all of these together into a single coherent, low-latency operational picture is the central technical challenge of sensor fusion for border security. India's diverse terrain makes this harder still, as the conditions across border sectors range from arid desert to dense riverine mangrove, and connectivity through radio and satellite links is frequently intermittent, adding further complexity to any integration effort.

\subsubsection{Existing Approaches}
The most widely referenced framework for structuring this problem is the Joint Directors of Laboratories (JDL) data fusion model, which defines five processing levels ranging from raw signal refinement at Level~0 through situation assessment at Level~2 and up to threat refinement at Level~3. It remains the dominant architectural reference for military fusion systems worldwide~\cite{meng2023multisource}. A comprehensive survey of multi-source information fusion for command, control, communications, and intelligence systems by Meng et al.~\cite{meng2023multisource} draws a useful distinction between data-level, feature-level, and decision-level fusion, mapping out the computational tradeoffs involved at each stage. In counter-UAV applications where individual sensor reliability varies across optical, thermal, radio-frequency, and radar inputs, attention-based cross-modal feature alignment has consistently produced the strongest classification accuracy. At the industry level, defence contractors have pushed this further by integrating hyperspectral imaging with AI-enabled radar and LiDAR processing for military situational awareness, and passive infrared sensor networks fused with multi-modal inputs have already been deployed for autonomous detect-classify-track operations at military installations. Ali et al.~\cite{threatdetect2024rt} survey transformer-based models and federated learning pipelines for real-time threat detection, identifying attention mechanisms as particularly effective for aligning features across sensor modalities.

\subsubsection{Key Insight and Research Gap}
Despite this progress, a consistent limitation runs through the reviewed literature. Nearly all fusion architectures have been evaluated on homogeneous or pairwise sensor combinations and have not been tested against the five-or-more modality configurations that are actually present in deployed systems like India's Comprehensive Integrated Border Management System (CIBMS). The temporal misalignment between high-frequency ground sensors reporting at 100~Hz or above and low-frequency satellite revisit passes separated by hours remains a genuinely unsolved synchronisation problem~\cite{meng2023multisource}. Perhaps most critically, not a single study in the reviewed corpus evaluates fusion performance under adversarial sensor degradation conditions. Given the documented electromagnetic interference threats present in contested border environments~\cite{milai2025risk}, this is not a minor academic gap but a direct operational vulnerability that future work urgently needs to address.


\subsection{Multi-Agent AI for Threat Detection and Classification}
\label{sec:multiagent}

\subsubsection{Problem Definition}
A single AI model simply cannot do everything at once. Asking one model to simultaneously process live video from dozens of cameras, classify radar anomalies, analyse seismic and acoustic time-series data, and correlate all of it into a unified threat picture within the sub-second response windows that operational protocols demand is not realistic~\cite{zhao2022deepai}. The natural solution is to distribute these responsibilities across multiple specialised agents, each focused on what it does best. But this introduces its own set of challenges. These agents need to share findings with each other, reconcile situations where two agents reach contradictory conclusions about the same event, and ultimately agree on a single threat classification without any one agent becoming a bottleneck that slows the entire system down~\cite{hassan2025multiagent}.

\subsubsection{Existing Approaches}
Several recent systems have explored how to structure this kind of multi-agent design effectively. Hassan et al.~\cite{hassan2025multiagent} propose a modular multi-agent system for cybersecurity threat detection where separate agents handle log analysis, IP analysis, and threat correlation independently before feeding into a central correlation engine. Their results show that this role-based decomposition improves both detection accuracy and the ability to trace why a particular classification was made, compared to a single monolithic model. Zhao et al.~\cite{zhao2022deepai} describe a cooperative AI staff system managing over 500{,}000 simulated battlefield agents, delivering command-level situational awareness through automated spatial and temporal pattern correlation across a large operational area. The AgenticCyber framework~\cite{agenticcyber2025} takes multi-agent orchestration further by fusing cloud logs, surveillance video, and audio streams across modalities, reporting a 96.2\% F1 score at 420~milliseconds end-to-end latency. On the visual detection side, YOLOv8-family architectures have been successfully adapted for military targets in satellite imagery~\cite{satellite2024yolov8} and weapon identification in live surveillance feeds~\cite{weapon2023yolov8}. Project Maven's real-world deployment offered perhaps the most striking validation of the approach, with a 20-person AI-assisted cell matching the video analysis throughput of a 2{,}000-person manual team~\cite{cset2024maven}, though independent evaluation found tank identification accuracy was only around 60\% even under favourable conditions~\cite{rand2023maven}. Frameworks for formally dividing decision authority between AI agents and human commanders in tactical settings have also been developed~\cite{autonomyteaming2024}.

\subsubsection{Key Insight and Research Gap}
Looking across these systems, a clear and important gap emerges. Existing multi-agent surveillance architectures have largely been designed around uniform task distributions, such as passing a tracked target between cameras in a consistent video network. None of them seriously tackle the problem of orchestrating agents that are reasoning over fundamentally different types of data at the same time, where one agent is working with pixel arrays, another with audio spectrograms, and another with radar range-Doppler matrices~\cite{hassan2025multiagent, agenticcyber2025}. This matters enormously for a system like VIGILANCE. A graduated five-level threat classification scheme requires every individual agent to produce well-calibrated confidence scores, and the correlation agent must then perform principled uncertainty quantification when combining outputs from agents that each have their own distinct error profiles and failure modes. In current multi-agent designs, confidence calibration is treated as a secondary concern rather than something built into the architecture from the ground up~\cite{cset2025ai}. That is a gap this work directly sets out to address.


\subsection{False-Alarm Reduction and Explainable Decision Support}
\label{sec:falsealarm}

\subsubsection{Problem Definition}
Persistently high false-alarm rates constitute a documented and operationally consequential failure mode in fielded surveillance systems~\cite{deepsurvey2023anomaly}. Sustained spurious alerts induce operator alert fatigue, misallocate rapid-response resources, and progressively degrade trust in automated detection pipelines. Root causes include environmental noise from wind-induced fence vibration, animal crossings, and weather-driven radar clutter; absence of temporal consistency verification in threshold-based alarm logic; and lack of contextual priors that distinguish known civilian activity zones from restricted perimeters~\cite{cctv2024anomaly}.

\subsubsection{Existing Approaches}
Weakly-supervised multiple-instance learning (MIL) frameworks for anomaly detection in real-world surveillance video have established training paradigms that substantially reduce annotation requirements for large-scale deployment~\cite{deepsurvey2023anomaly}. Deep-learning surveys catalogue CNN, LSTM, and transformer-based temporal models that improve anomaly discrimination by enforcing signature persistence across consecutive frames before triggering an alert. Hybrid architectures combining spatial object detection with temporal verification have been shown to reduce false-positive rates relative to single-frame classifiers in CCTV monitoring~\cite{cctv2024anomaly}. On the explainability front, Zhang et al.~\cite{zhang2022xai} survey SHAP, LIME, and attention-based explanation mechanisms applicable to black-box classifiers, applying explainable AI (XAI) to intrusion detection and malware classification, tasks that are structurally analogous to the threat-triage function required in military surveillance. Roff and Danks~\cite{military2024xai} address military-specific accountability requirements, stipulating that commanders must be able to audit the rationale behind targeting and escalatory recommendations under both domestic and international legal obligations.

\subsubsection{Key Insight and Research Gap}
XAI methods in surveillance have been evaluated primarily for post-hoc interpretability on static or pre-recorded datasets~\cite{zhang2022xai} and have not been integrated into real-time decision pipelines where explanation generation latency itself constrains operator response time. Contextual filtering approaches that suppress alarms in geofenced civilian zones depend on accurate, current civilian activity maps that are unavailable or classified in active border areas~\cite{cctv2024anomaly}, creating an unresolved data-dependency bottleneck that no reviewed study has addressed.


\subsection{Geospatial Situational Awareness and Human-in-the-Loop Systems}
\label{sec:geospatial}

\subsubsection{Problem Definition}
Actionable command decisions require threat intelligence to be spatially anchored: operators must perceive not only the nature and severity of a detected anomaly but also its precise location relative to terrain features, installation boundaries, sensor coverage envelopes, and proximity to civilian population centres~\cite{karimipour2021iot}. Graduated threat classification further demands that alerts assessed at the highest severity levels do not trigger autonomous system action; human authorisation must remain mandatory at those tiers~\cite{meaningful2024hitl}. The interface design challenge is twofold: preventing operators from reflexively endorsing AI recommendations while preserving response latencies compatible with tactical decision cycles~\cite{hitl2023trust}.

\subsubsection{Existing Approaches}
Karimipour and Derakhshan~\cite{karimipour2021iot} describe AI-enabled situational awareness dashboards for distributed sensor networks in industrial Internet of Things (IoT) environments, incorporating geographic information system (GIS) threat mapping with sensor registries that visualise coverage zones, detection radii, and active alarm sources on unified terrain layers. Zhao et al.~\cite{zhao2022deepai} extend this concept to military command by constructing multi-layer spatio-temporal battlefield visualisations that overlay force positions, predicted adversary trajectories, and terrain accessibility indices. Geospatial UAV detection with movement-vector plotting on georeferenced maps has been demonstrated for counter-drone applications. Maas~\cite{hitl2023trust} examines trust calibration and automation bias in military human-autonomy teaming, reporting that operator trust is sensitive to the transparency and perceived reliability of AI outputs. Svenmarck and Luotsinen~\cite{meaningful2024hitl} formalise the legal and technical requirements for meaningful human control, proposing design constraints that ensure commander authority over engagement decisions is not diluted by automated recommendation systems.

\subsubsection{Key Insight and Research Gap}
Human-in-the-loop (HITL) frameworks in the surveyed literature focus on single-operator, single-threat interaction models~\cite{hitl2023trust, meaningful2024hitl}. Concurrent multi-threat environments, in which an operator must triage AI-generated alerts across geographically distributed incidents simultaneously, remain under-examined despite constituting the operational norm in border-surveillance command centres. Interface designs that communicate prediction uncertainty quantitatively to operators with varying levels of technical expertise present an open challenge; existing prototypes either display raw probability values that non-technical operators misinterpret or suppress uncertainty information entirely, promoting uncalibrated trust~\cite{hitl2023trust}.


\subsection{Real-Time Distributed Architecture and Adversarial Robustness}
\label{sec:distributed}

\subsubsection{Problem Definition}
Sensor data generated at a remote border outpost must traverse constrained communication links, undergo inference at sub-second latency, and arrive at a central command dashboard with actionable fidelity~\cite{flink2024kafka}. Meeting this requirement demands edge preprocessing to reduce transmission volume, stream-processing middleware for sequential low-latency inference, and an architecture tolerant of node failures and intermittent connectivity. Adversarial robustness introduces a distinct threat surface: a high-value surveillance system constitutes an explicit target for adversarial patches designed to defeat object detectors, GPS spoofing aimed at unmanned platform navigation, and signal jamming directed at ground-sensor communication uplinks~\cite{advyolo2022, milai2025risk}.

\subsubsection{Existing Approaches}
Patel et al.~\cite{flink2024kafka} present a real-time IoT sensor streaming architecture combining Apache Kafka for telemetry ingestion with Apache Flink for stream processing, achieving low-latency anomaly detection on continuous sensor feeds with fault-tolerant horizontal scaling. Edge-computing extensions of this pattern perform preliminary inference at the sensor node before transmitting compressed results, reducing bandwidth requirements for drone and CCTV data in remote deployments. Ali et al.~\cite{threatdetect2024rt} survey transformer-based models and federated learning pipelines for real-time threat detection across IoT and 5G networks, identifying federated architectures as particularly suitable for privacy-preserving distributed inference in military contexts. On the adversarial front, Liang et al.~\cite{advyolo2022} catalogue physical-world attacks on YOLO-family detectors, including adversarial patches and printed perturbations that cause vehicle misclassification in outdoor conditions. Wang et al.~\cite{milai2025risk} address military-specific adversarial risks such as electromagnetic interference and GPS spoofing, and propose ethical governance guardrails for deployed military AI. Scientometric analysis of approximately 3{,}500 Scopus records confirms a fourfold growth in AI-related defence publications between 2015 and 2025~\cite{emerging2025scientometric}, yet distributed surveillance architectures remain under-represented relative to single-algorithm detection studies. Independent PRISMA reviews of UAV cybersecurity~\cite{uavsec2025iet} and UAV intrusion detection systems~\cite{uavids2025drones} converge on a shared conclusion: no unified cross-domain framework currently addresses the full defence surveillance pipeline from edge to command centre.

\subsubsection{Key Insight and Research Gap}
Adversarial robustness research in computer vision predominantly addresses digital perturbations on laboratory-curated datasets; physical-world attacks under outdoor surveillance conditions involving variable lighting, partial occlusion, and ambient sensor noise are substantially less investigated~\cite{advyolo2022}. No reviewed distributed-architecture study evaluates end-to-end pipeline latency under simultaneous edge-node failure and adversarial sensor injection, the compound threat model most representative of a contested border environment. The gap between adversarial machine-learning research and operational distributed systems engineering constitutes one of the least-addressed intersections in the defence surveillance literature~\cite{uavsec2025iet, milai2025risk}.


\section{Comparative Analysis}
\label{sec:comparative}

Cross-domain evaluation reveals an uneven maturity gradient across the five surveyed areas. Sensor fusion (Domain~A) and visual object detection (a component of Domain~B) represent the most technically mature pillars, and YOLOv8-family detectors~\cite{satellite2024yolov8, weapon2023yolov8} approach deployment readiness for constrained single-modality scenarios. False-alarm reduction techniques (Domain~C) yield measurable improvements on benchmark datasets~\cite{deepsurvey2023anomaly} but lack validated integration with real-time explainability pipelines and with the accountability demands specific to military decision chains~\cite{military2024xai}. Geospatial HITL systems (Domain~D) and adversarially hardened distributed architectures (Domain~E) remain at the proof-of-concept stage, with empirical validation limited to single-site deployments or simulated environments~\cite{karimipour2021iot, flink2024kafka}.

The most consequential pattern emerging from this review is the near-total absence of research addressing interaction effects across domain boundaries. Fusion fidelity at the sensor-integration layer directly determines the confidence bounds available to downstream classification agents~\cite{meng2023multisource, hassan2025multiagent}. Explanation-generation overhead in the XAI layer imposes an additive latency penalty on the HITL authorisation cycle~\cite{zhang2022xai, hitl2023trust}. Adversarial corruption at edge sensors propagates cascading failures through fusion, classification, and alerting stages that no surveyed system models as an integrated failure mode~\cite{advyolo2022, milai2025risk}. Table~\ref{tab:comparative} summarises the best-performing technique, reported benchmark, metric, and principal limitation for each domain when assessed against the requirements of a unified multi-domain surveillance framework.

\begin{table*}[!t]
\centering
\caption{Comparative Analysis of Surveyed Approaches}
\label{tab:comparative}
\renewcommand{\arraystretch}{1.3}
\footnotesize
\begin{tabular}{>{\raggedright\arraybackslash}p{1.5cm}>{\raggedright\arraybackslash}p{3.2cm}>{\raggedright\arraybackslash}p{3.0cm}>{\raggedright\arraybackslash}p{2.8cm}>{\raggedright\arraybackslash}p{3.5cm}}
\hline
\textbf{Domain} & \textbf{Best Performing Technique} & \textbf{Dataset / Benchmark} & \textbf{Metric Reported} & \textbf{Limitation in Unified Context} \\
\hline
A: Sensor Fusion
  & Attention-based cross-modal DL fusion
  & Counter-UAV optical, thermal, RF, radar corpus
  & Highest classification accuracy vs.\ single-modal baselines
  & No temporal synchronisation for $>$2 modalities at disparate sampling rates \\
\hline
B: Multi-Agent AI
  & Multi-modal agentic orchestration
  & Cloud logs + surveillance video + audio
  & 96.2\% F1; 420\,ms end-to-end latency
  & Agent confidence calibration treated as post-hoc addition \\
\hline
C: False Alarm / XAI
  & Deep-learning anomaly detection with temporal verification
  & Real-world surveillance video benchmarks
  & Benchmark AUC across anomaly classes
  & Explanation latency not measured for real-time deployment \\
\hline
D: Geospatial / HITL
  & Multi-layer battlefield visualisation
  & 500K+ simulated battlefield agents
  & Command-level situational awareness coverage
  & Multi-threat concurrent operator triage not evaluated \\
\hline
E: Distributed / Adversarial
  & Kafka + Flink real-time streaming
  & Continuous IoT sensor streams
  & Low-latency, fault-tolerant stream processing
  & Not tested under adversarial sensor injection or node compromise \\
\hline
\end{tabular}
\end{table*}
 from main.tex
% ============================================================

\section{Taxonomy of Related Work}
\label{sec:taxonomy}

The surveyed literature was organized into five functional domains that together cover the full processing pipeline of an integrated defence surveillance system: (A) multi-modal sensor fusion and data integration, (B) multi-agent AI for threat detection and classification, (C) false-alarm reduction and explainable decision support, (D) geospatial situational awareness and human-in-the-loop governance, and (E) real-time distributed architecture and adversarial robustness. To keep the scope focused and relevant, inclusion criteria limited the corpus to studies that satisfied at least two of three conditions: applicability to defence or security-critical environments, use of AI or machine learning methods, and involvement of multi-sensor or distributed system architectures. The sources drawn upon span peer-reviewed journal articles, IEEE and ACM conference proceedings, and policy-research technical reports~\cite{cset2024maven, cset2025ai}.

Table~\ref{tab:taxonomy} presents the resulting classification. For each domain, the table identifies representative published approaches, the principal sensor or data modalities examined, the key technique employed, and the primary limitation observed when that approach is evaluated against the demands of a unified multi-domain surveillance framework. It is worth noting that these five domains do not operate in isolation from one another in practice. The quality of sensor fusion at the data integration layer in Domain A directly shapes the classification confidence that downstream agents in Domain B can reliably achieve. Similarly, the time it takes to generate explanations in the decision-support layer of Domain C places real constraints on how quickly human commanders can complete their authorization loop in Domain D. These cross-domain dependencies are examined quantitatively in Section~V.



\section{Literature Survey}
\label{sec:survey}

\subsection{Multi-Modal Sensor Fusion and Data Integration}
\label{sec:fusion}

\subsubsection{Problem Definition}
One of the most fundamental challenges in defence surveillance is that the sensors deployed across a border do not speak the same language. Satellite imagery covers vast areas but only updates every few hours. UAV feeds capture narrow zones in high-frame-rate video. Ground-based sensors like seismic detectors, acoustic microphones, and infrared cameras produce continuous streams of numeric time-series data, each with its own noise characteristics and sampling rate~\cite{meng2023multisource}. Bringing all of these together into a single coherent, low-latency operational picture is the central technical challenge of sensor fusion for border security. India's diverse terrain makes this harder still, as the conditions across border sectors range from arid desert to dense riverine mangrove, and connectivity through radio and satellite links is frequently intermittent, adding further complexity to any integration effort.

\subsubsection{Existing Approaches}
The most widely referenced framework for structuring this problem is the Joint Directors of Laboratories (JDL) data fusion model, which defines five processing levels ranging from raw signal refinement at Level~0 through situation assessment at Level~2 and up to threat refinement at Level~3. It remains the dominant architectural reference for military fusion systems worldwide~\cite{meng2023multisource}. A comprehensive survey of multi-source information fusion for command, control, communications, and intelligence systems by Meng et al.~\cite{meng2023multisource} draws a useful distinction between data-level, feature-level, and decision-level fusion, mapping out the computational tradeoffs involved at each stage. In counter-UAV applications where individual sensor reliability varies across optical, thermal, radio-frequency, and radar inputs, attention-based cross-modal feature alignment has consistently produced the strongest classification accuracy. At the industry level, defence contractors have pushed this further by integrating hyperspectral imaging with AI-enabled radar and LiDAR processing for military situational awareness, and passive infrared sensor networks fused with multi-modal inputs have already been deployed for autonomous detect-classify-track operations at military installations. Ali et al.~\cite{threatdetect2024rt} survey transformer-based models and federated learning pipelines for real-time threat detection, identifying attention mechanisms as particularly effective for aligning features across sensor modalities.

\subsubsection{Key Insight and Research Gap}
Despite this progress, a consistent limitation runs through the reviewed literature. Nearly all fusion architectures have been evaluated on homogeneous or pairwise sensor combinations and have not been tested against the five-or-more modality configurations that are actually present in deployed systems like India's Comprehensive Integrated Border Management System (CIBMS). The temporal misalignment between high-frequency ground sensors reporting at 100~Hz or above and low-frequency satellite revisit passes separated by hours remains a genuinely unsolved synchronisation problem~\cite{meng2023multisource}. Perhaps most critically, not a single study in the reviewed corpus evaluates fusion performance under adversarial sensor degradation conditions. Given the documented electromagnetic interference threats present in contested border environments~\cite{milai2025risk}, this is not a minor academic gap but a direct operational vulnerability that future work urgently needs to address.


\subsection{Multi-Agent AI for Threat Detection and Classification}
\label{sec:multiagent}

\subsubsection{Problem Definition}
A single AI model simply cannot do everything at once. Asking one model to simultaneously process live video from dozens of cameras, classify radar anomalies, analyse seismic and acoustic time-series data, and correlate all of it into a unified threat picture within the sub-second response windows that operational protocols demand is not realistic~\cite{zhao2022deepai}. The natural solution is to distribute these responsibilities across multiple specialised agents, each focused on what it does best. But this introduces its own set of challenges. These agents need to share findings with each other, reconcile situations where two agents reach contradictory conclusions about the same event, and ultimately agree on a single threat classification without any one agent becoming a bottleneck that slows the entire system down~\cite{hassan2025multiagent}.

\subsubsection{Existing Approaches}
Several recent systems have explored how to structure this kind of multi-agent design effectively. Hassan et al.~\cite{hassan2025multiagent} propose a modular multi-agent system for cybersecurity threat detection where separate agents handle log analysis, IP analysis, and threat correlation independently before feeding into a central correlation engine. Their results show that this role-based decomposition improves both detection accuracy and the ability to trace why a particular classification was made, compared to a single monolithic model. Zhao et al.~\cite{zhao2022deepai} describe a cooperative AI staff system managing over 500{,}000 simulated battlefield agents, delivering command-level situational awareness through automated spatial and temporal pattern correlation across a large operational area. The AgenticCyber framework~\cite{agenticcyber2025} takes multi-agent orchestration further by fusing cloud logs, surveillance video, and audio streams across modalities, reporting a 96.2\% F1 score at 420~milliseconds end-to-end latency. On the visual detection side, YOLOv8-family architectures have been successfully adapted for military targets in satellite imagery~\cite{satellite2024yolov8} and weapon identification in live surveillance feeds~\cite{weapon2023yolov8}. Project Maven's real-world deployment offered perhaps the most striking validation of the approach, with a 20-person AI-assisted cell matching the video analysis throughput of a 2{,}000-person manual team~\cite{cset2024maven}, though independent evaluation found tank identification accuracy was only around 60\% even under favourable conditions~\cite{rand2023maven}. Frameworks for formally dividing decision authority between AI agents and human commanders in tactical settings have also been developed~\cite{autonomyteaming2024}.

\subsubsection{Key Insight and Research Gap}
Looking across these systems, a clear and important gap emerges. Existing multi-agent surveillance architectures have largely been designed around uniform task distributions, such as passing a tracked target between cameras in a consistent video network. None of them seriously tackle the problem of orchestrating agents that are reasoning over fundamentally different types of data at the same time, where one agent is working with pixel arrays, another with audio spectrograms, and another with radar range-Doppler matrices~\cite{hassan2025multiagent, agenticcyber2025}. This matters enormously for a system like VIGILANCE. A graduated five-level threat classification scheme requires every individual agent to produce well-calibrated confidence scores, and the correlation agent must then perform principled uncertainty quantification when combining outputs from agents that each have their own distinct error profiles and failure modes. In current multi-agent designs, confidence calibration is treated as a secondary concern rather than something built into the architecture from the ground up~\cite{cset2025ai}. That is a gap this work directly sets out to address.


\subsection{False-Alarm Reduction and Explainable Decision Support}
\label{sec:falsealarm}

\subsubsection{Problem Definition}
Persistently high false-alarm rates constitute a documented and operationally consequential failure mode in fielded surveillance systems~\cite{deepsurvey2023anomaly}. Sustained spurious alerts induce operator alert fatigue, misallocate rapid-response resources, and progressively degrade trust in automated detection pipelines. Root causes include environmental noise from wind-induced fence vibration, animal crossings, and weather-driven radar clutter; absence of temporal consistency verification in threshold-based alarm logic; and lack of contextual priors that distinguish known civilian activity zones from restricted perimeters~\cite{cctv2024anomaly}.

\subsubsection{Existing Approaches}
Weakly-supervised multiple-instance learning (MIL) frameworks for anomaly detection in real-world surveillance video have established training paradigms that substantially reduce annotation requirements for large-scale deployment~\cite{deepsurvey2023anomaly}. Deep-learning surveys catalogue CNN, LSTM, and transformer-based temporal models that improve anomaly discrimination by enforcing signature persistence across consecutive frames before triggering an alert. Hybrid architectures combining spatial object detection with temporal verification have been shown to reduce false-positive rates relative to single-frame classifiers in CCTV monitoring~\cite{cctv2024anomaly}. On the explainability front, Zhang et al.~\cite{zhang2022xai} survey SHAP, LIME, and attention-based explanation mechanisms applicable to black-box classifiers, applying explainable AI (XAI) to intrusion detection and malware classification, tasks that are structurally analogous to the threat-triage function required in military surveillance. Roff and Danks~\cite{military2024xai} address military-specific accountability requirements, stipulating that commanders must be able to audit the rationale behind targeting and escalatory recommendations under both domestic and international legal obligations.

\subsubsection{Key Insight and Research Gap}
XAI methods in surveillance have been evaluated primarily for post-hoc interpretability on static or pre-recorded datasets~\cite{zhang2022xai} and have not been integrated into real-time decision pipelines where explanation generation latency itself constrains operator response time. Contextual filtering approaches that suppress alarms in geofenced civilian zones depend on accurate, current civilian activity maps that are unavailable or classified in active border areas~\cite{cctv2024anomaly}, creating an unresolved data-dependency bottleneck that no reviewed study has addressed.


\subsection{Geospatial Situational Awareness and Human-in-the-Loop Systems}
\label{sec:geospatial}

\subsubsection{Problem Definition}
Actionable command decisions require threat intelligence to be spatially anchored: operators must perceive not only the nature and severity of a detected anomaly but also its precise location relative to terrain features, installation boundaries, sensor coverage envelopes, and proximity to civilian population centres~\cite{karimipour2021iot}. Graduated threat classification further demands that alerts assessed at the highest severity levels do not trigger autonomous system action; human authorisation must remain mandatory at those tiers~\cite{meaningful2024hitl}. The interface design challenge is twofold: preventing operators from reflexively endorsing AI recommendations while preserving response latencies compatible with tactical decision cycles~\cite{hitl2023trust}.

\subsubsection{Existing Approaches}
Karimipour and Derakhshan~\cite{karimipour2021iot} describe AI-enabled situational awareness dashboards for distributed sensor networks in industrial Internet of Things (IoT) environments, incorporating geographic information system (GIS) threat mapping with sensor registries that visualise coverage zones, detection radii, and active alarm sources on unified terrain layers. Zhao et al.~\cite{zhao2022deepai} extend this concept to military command by constructing multi-layer spatio-temporal battlefield visualisations that overlay force positions, predicted adversary trajectories, and terrain accessibility indices. Geospatial UAV detection with movement-vector plotting on georeferenced maps has been demonstrated for counter-drone applications. Maas~\cite{hitl2023trust} examines trust calibration and automation bias in military human-autonomy teaming, reporting that operator trust is sensitive to the transparency and perceived reliability of AI outputs. Svenmarck and Luotsinen~\cite{meaningful2024hitl} formalise the legal and technical requirements for meaningful human control, proposing design constraints that ensure commander authority over engagement decisions is not diluted by automated recommendation systems.

\subsubsection{Key Insight and Research Gap}
Human-in-the-loop (HITL) frameworks in the surveyed literature focus on single-operator, single-threat interaction models~\cite{hitl2023trust, meaningful2024hitl}. Concurrent multi-threat environments, in which an operator must triage AI-generated alerts across geographically distributed incidents simultaneously, remain under-examined despite constituting the operational norm in border-surveillance command centres. Interface designs that communicate prediction uncertainty quantitatively to operators with varying levels of technical expertise present an open challenge; existing prototypes either display raw probability values that non-technical operators misinterpret or suppress uncertainty information entirely, promoting uncalibrated trust~\cite{hitl2023trust}.


\subsection{Real-Time Distributed Architecture and Adversarial Robustness}
\label{sec:distributed}

\subsubsection{Problem Definition}
Sensor data generated at a remote border outpost must traverse constrained communication links, undergo inference at sub-second latency, and arrive at a central command dashboard with actionable fidelity~\cite{flink2024kafka}. Meeting this requirement demands edge preprocessing to reduce transmission volume, stream-processing middleware for sequential low-latency inference, and an architecture tolerant of node failures and intermittent connectivity. Adversarial robustness introduces a distinct threat surface: a high-value surveillance system constitutes an explicit target for adversarial patches designed to defeat object detectors, GPS spoofing aimed at unmanned platform navigation, and signal jamming directed at ground-sensor communication uplinks~\cite{advyolo2022, milai2025risk}.

\subsubsection{Existing Approaches}
Patel et al.~\cite{flink2024kafka} present a real-time IoT sensor streaming architecture combining Apache Kafka for telemetry ingestion with Apache Flink for stream processing, achieving low-latency anomaly detection on continuous sensor feeds with fault-tolerant horizontal scaling. Edge-computing extensions of this pattern perform preliminary inference at the sensor node before transmitting compressed results, reducing bandwidth requirements for drone and CCTV data in remote deployments. Ali et al.~\cite{threatdetect2024rt} survey transformer-based models and federated learning pipelines for real-time threat detection across IoT and 5G networks, identifying federated architectures as particularly suitable for privacy-preserving distributed inference in military contexts. On the adversarial front, Liang et al.~\cite{advyolo2022} catalogue physical-world attacks on YOLO-family detectors, including adversarial patches and printed perturbations that cause vehicle misclassification in outdoor conditions. Wang et al.~\cite{milai2025risk} address military-specific adversarial risks such as electromagnetic interference and GPS spoofing, and propose ethical governance guardrails for deployed military AI. Scientometric analysis of approximately 3{,}500 Scopus records confirms a fourfold growth in AI-related defence publications between 2015 and 2025~\cite{emerging2025scientometric}, yet distributed surveillance architectures remain under-represented relative to single-algorithm detection studies. Independent PRISMA reviews of UAV cybersecurity~\cite{uavsec2025iet} and UAV intrusion detection systems~\cite{uavids2025drones} converge on a shared conclusion: no unified cross-domain framework currently addresses the full defence surveillance pipeline from edge to command centre.

\subsubsection{Key Insight and Research Gap}
Adversarial robustness research in computer vision predominantly addresses digital perturbations on laboratory-curated datasets; physical-world attacks under outdoor surveillance conditions involving variable lighting, partial occlusion, and ambient sensor noise are substantially less investigated~\cite{advyolo2022}. No reviewed distributed-architecture study evaluates end-to-end pipeline latency under simultaneous edge-node failure and adversarial sensor injection, the compound threat model most representative of a contested border environment. The gap between adversarial machine-learning research and operational distributed systems engineering constitutes one of the least-addressed intersections in the defence surveillance literature~\cite{uavsec2025iet, milai2025risk}.


\section{Comparative Analysis}
\label{sec:comparative}

Cross-domain evaluation reveals an uneven maturity gradient across the five surveyed areas. Sensor fusion (Domain~A) and visual object detection (a component of Domain~B) represent the most technically mature pillars, and YOLOv8-family detectors~\cite{satellite2024yolov8, weapon2023yolov8} approach deployment readiness for constrained single-modality scenarios. False-alarm reduction techniques (Domain~C) yield measurable improvements on benchmark datasets~\cite{deepsurvey2023anomaly} but lack validated integration with real-time explainability pipelines and with the accountability demands specific to military decision chains~\cite{military2024xai}. Geospatial HITL systems (Domain~D) and adversarially hardened distributed architectures (Domain~E) remain at the proof-of-concept stage, with empirical validation limited to single-site deployments or simulated environments~\cite{karimipour2021iot, flink2024kafka}.

The most consequential pattern emerging from this review is the near-total absence of research addressing interaction effects across domain boundaries. Fusion fidelity at the sensor-integration layer directly determines the confidence bounds available to downstream classification agents~\cite{meng2023multisource, hassan2025multiagent}. Explanation-generation overhead in the XAI layer imposes an additive latency penalty on the HITL authorisation cycle~\cite{zhang2022xai, hitl2023trust}. Adversarial corruption at edge sensors propagates cascading failures through fusion, classification, and alerting stages that no surveyed system models as an integrated failure mode~\cite{advyolo2022, milai2025risk}. Table~\ref{tab:comparative} summarises the best-performing technique, reported benchmark, metric, and principal limitation for each domain when assessed against the requirements of a unified multi-domain surveillance framework.

\begin{table*}[!t]
\centering
\caption{Comparative Analysis of Surveyed Approaches}
\label{tab:comparative}
\renewcommand{\arraystretch}{1.3}
\footnotesize
\begin{tabular}{>{\raggedright\arraybackslash}p{1.5cm}>{\raggedright\arraybackslash}p{3.2cm}>{\raggedright\arraybackslash}p{3.0cm}>{\raggedright\arraybackslash}p{2.8cm}>{\raggedright\arraybackslash}p{3.5cm}}
\hline
\textbf{Domain} & \textbf{Best Performing Technique} & \textbf{Dataset / Benchmark} & \textbf{Metric Reported} & \textbf{Limitation in Unified Context} \\
\hline
A: Sensor Fusion
  & Attention-based cross-modal DL fusion
  & Counter-UAV optical, thermal, RF, radar corpus
  & Highest classification accuracy vs.\ single-modal baselines
  & No temporal synchronisation for $>$2 modalities at disparate sampling rates \\
\hline
B: Multi-Agent AI
  & Multi-modal agentic orchestration
  & Cloud logs + surveillance video + audio
  & 96.2\% F1; 420\,ms end-to-end latency
  & Agent confidence calibration treated as post-hoc addition \\
\hline
C: False Alarm / XAI
  & Deep-learning anomaly detection with temporal verification
  & Real-world surveillance video benchmarks
  & Benchmark AUC across anomaly classes
  & Explanation latency not measured for real-time deployment \\
\hline
D: Geospatial / HITL
  & Multi-layer battlefield visualisation
  & 500K+ simulated battlefield agents
  & Command-level situational awareness coverage
  & Multi-threat concurrent operator triage not evaluated \\
\hline
E: Distributed / Adversarial
  & Kafka + Flink real-time streaming
  & Continuous IoT sensor streams
  & Low-latency, fault-tolerant stream processing
  & Not tested under adversarial sensor injection or node compromise \\
\hline
\end{tabular}
\end{table*}
